\scanpage{67}
\mspage{29}
Enfin, la condition d) signifie que cette structure
\inserted{d'espace top. ordonné}
(satisfait la condition c) du Scholie (p. $27$).

\uline{Cas équisingulier} : les $s_{ij}$ sont fibrants propres

\uline{Cas non singulier} : les $\Sigma_i$ sont des variétés
topologiques; le bord de $\Sigma_i$ est la réunion
des $\Sigma_{ji}$ pour $j<i$ \deleted{(NB il suffit
de prendre pour $j$}
\margin{On voudra que les $\Sigma_{ij}$ soient des variétés à bord, que le
bord soit la réunion des $\Sigma_{kij}$ et $\Sigma_{ikj}$ --- mais
c'est sans doute une conséquence de la condition précédente\ldots}

\uline{Cas équisingulier-non singulier}
$s_{ij}:\Sigma_{ik}\longrightarrow\unsure{\Sigma_i|\Sigma_j}$ est
une fibration propre, à fibres des variétés
à bord, le fibré des bords est réunion
des $\Sigma_{ijk}$ pour $i<j<k$.

13. \uline{Une digression} Revenons au cas de
$\Sigma,\unsure{\Sigma^{*}(1)|\Sigma^{<}(1)}\subset\Sigma\times\Sigma$, sans donnée de $I$. Plus
généralement, considérons une \uline{catégorie topologique}
\[
\uline{\Sigma}=(\Sigma_{(0)}=\operatorname{Ob}(\uline{\Sigma}),\quad\Sigma_{(1)}=\operatorname{Fl}(\uline{\Sigma}),
\]
\diagram{\begin{tikzcd}
\Sigma_{(1)}\arrow[r,shift left=2,"\uline{s}"]\arrow[r,shift right=2,"\uline{b}"']&\Sigma_{(0)}
\end{tikzcd}}
\diagram{\begin{tikzcd}[column sep=large]
(\Sigma_{(1)},b)\times_{\Sigma_{(0)}}(\Sigma_{(1)},s)\arrow[r,Rightarrow,"\text{déf}"] & \Sigma_{(2)}\arrow[r,"\unsure{\text{cpl.}|\text{compl.}|\text{appl.}}","\text{comp. des flèches}"'] & \Sigma_{(1)}
\end{tikzcd}}
), et supposons
que $\uline{s}:\Sigma_{(1)}\longrightarrow\Sigma_{(0)}$ soit \uline{fibrant}. Je dis
qu'alors $\pi_0(\Sigma_{(0)})=\unsure{I_{(0)}|I}$ est \uline{un ens. \inserted{pré}ordonné},
en prenant comme graphe de la relation
d'ordre l'image de
\diagram{\begin{tikzcd}[column sep=large]
\pi_0(\Sigma_{(1)})\arrow[r,"{\pi_0(\uline{s},\uline{b})}"]&\pi_0(\Sigma_{(0)}\times\Sigma_{(0)})\arrow[d,"\sim"]\\
&\pi_0(\Sigma_{(0)}\times\Sigma_{(0)})
\end{tikzcd}}
\mspage{66}
